\documentclass[11pt]{article}
\usepackage[T1]{fontenc}
\usepackage{textcomp}
\usepackage[margin=0.75in]{geometry}
\usepackage{amsmath,amssymb,amsthm}
\usepackage{array,booktabs}
\usepackage{hyperref}
\setlength{\parindent}{0pt}
\newtheorem{theorem}{Theorem}
\theoremstyle{definition}
\newtheorem{definition}{Definition}
\title{Flow Proof Artifact: prop-12-apply-parallelogram-exceeding}
\author{Generated by \texttt{flow doc proof}}
\date{Flow Proof Book}
\begin{document}
\maketitle
To a given straight line to apply a parallelogram equal to a given rectilineal figure and exceeding by a parallelogram similar to a given one.\\[0.5em]
\textbf{Source.} euclid — Elements, Book VI, Proposition 12\\[0.5em]
\bigskip
\noindent{\large \textbf{Derived fact 1.} \textit{To a given straight line to apply a parallelogram equal to a given rectilineal figure and exceeding by a parallelogram similar to a given one} \hfill the Euclidean plane \cdot Euclid Book VI \cdot Proposition 12: to apply a parallelogram equal to a figure and exceeding by a similar parallelogram}\par
\smallskip\noindent\textit{Coordinate: the Euclidean plane \cdot Euclid Book VI \cdot Proposition 12: to apply a parallelogram equal to a figure and exceeding by a similar parallelogram}\par
\smallskip\noindent\textit{Source: euclid — Elements, Book VI, Proposition 12}\par
\smallskip\noindent\textit{Built on: proposition 11: to apply a parallelogram equal to a figure and deficient by a similar parallelogram, for Euclid Book VI on the Euclidean plane}\par
\medskip
\noindent\textbf{Goal.} To a given straight line to apply a parallelogram equal to a given rectilineal figure and exceeding by a parallelogram similar to a given one\par
\noindent$\displaystyle \text{exceeding parallelogram applied}$\par
\medskip
\noindent
\begin{tabular}{@{} >{\raggedright\arraybackslash}p{0.06\textwidth} >{\raggedright\arraybackslash}p{0.47\textwidth} >{\raggedleft\arraybackslash}p{0.41\textwidth} @{}}
\textbf{\#} & \textbf{Proof} & \textbf{Mathematics} \\
\midrule
\textbf{1.} & We prove that proposition 12: to apply a parallelogram equal to a figure and exceeding by a similar parallelogram for Euclid Book VI the Euclidean plane. & \\[0.45em]
\textbf{2.} & Let AB = given\_line. & \\[0.45em]
\textbf{3.} & Let F = given\_figure. & \\[0.45em]
\textbf{4.} & Let par = applied\_parallelogram. & \\[0.45em]
\textbf{5.} & From step 2, step 3, and step 4, we can deduce that area(par) equals area(F) plus excess. & $\displaystyle area(par) = area(F) + excess$ \\[0.45em]
\textbf{6.} & We invoke the derived fact governing Euclid Book VI the Euclidean plane: proposition 11: to apply a parallelogram equal to a figure and deficient by a similar parallelogram, for Euclid Book VI on the Euclidean plane. & \\[0.45em]
\textbf{7.} & From step 6, this implies exceeding parallelogram applied. Hence proven. & $\displaystyle \text{exceeding parallelogram applied}$ \\[0.45em]
\bottomrule
\end{tabular}
\label{thm:Geometry-Euclid Book VI-proposition_12_to_apply_a_parallelogram_equal_to_a_figure_and_exceeding_by_a_similar_parallelogram}
\par\medskip\hrule
\end{document}
